
\chapter{The removable uncertainty}

\textbf{Purpose:} Divide the CODATA listing into the uncertainties arithmetic can remove and the
ones it cannot, and carry out the removable ones. \textbf{Scope:}
\texttt{codata\_exact.py}, \texttt{dn\_constants.py}, \texttt{proved\_digits.py},
\texttt{reproduce.py}.

Three different types of number wear the word \emph{exact} in the 2022 CODATA listing. The result is
which is which.

\section{The split}

Counted from the listing as fetched from NIST on 2026-09-11:

\begin{center}
\begin{tabular}{lr}
\toprule
class & count \\
\midrule
constants in the listing & 358 \\
marked \texttt{(exact)} & 81 \\
\quad of those, \textbf{terminating} & 19 \\
\quad of those, \textbf{truncated} & \textbf{62} \\
carrying a measured uncertainty & \textbf{277} \\
\bottomrule
\end{tabular}
\end{center}

\textbf{The nineteen terminating ones need nothing.} \(1.602\,176\,634 \times 10^{-19}\) is the
elementary charge, and that is all. The printed digits are the value, with no tail.
A thousand places of them would only append 985 zeros.

\textbf{The sixty two truncated ones are the target.} NIST prints them as
\(1.054\,571\,817\ldots \times 10^{-34}\), and the ellipsis is not an uncertainty. It is a
typesetting decision. Every one is a closed form in the seven constants that define the SI, and those are
exact \emph{by definition} since the 2019 redefinition. Their absent digits are unprinted, and not
unknown.

\textbf{The two hundred and seventy seven measured ones are permanent.} The fine-structure constant,
the gravitational constant, the electron mass. Their uncertainty is what a laboratory measured.
Computing \(\alpha\) to a thousand places would be computing a thousand places of a number known to
eleven, and any claim that precision reduces an experimental error is false.

So the honest statement is not that the uncertainty was removed. It is \textbf{which uncertainties
are mathematical and removable, and which are physical and permanent}. That division is not
published anywhere and it falls straight out of the listing.

\section{What was carried out}

Forty four of the sixty two, at a thousand places.

\begin{center}
\begin{tabular}{lr}
\toprule
verification & count \\
\midrule
exactly verified by multiplying back & \textbf{38} \\
verified against this tree's own \(\pi\) & \textbf{6} \\
agreeing with the NIST printed digits & \textbf{44 of 44} \\
\bottomrule
\end{tabular}
\end{center}

The other eighteen need extra definitions, and not extra arithmetic: the conventional
\(-90\) values, the Loschmidt constants and molar volumes at stated gas conditions, the radiation
constants, the Stefan-Boltzmann constant which needs \(\pi^5\), and the two Wien displacement
constants, which need the root of \(x = 5(1 - e^{-x})\) and are the only entries here that are not a
closed form in the defining constants at all.

\section{Why this verification is stronger than the rest of the precision table}

\texttt{dn\_constants.csv} lives by two routes agreeing to a zero gap. That is evidence, and two
routes can share an error.

A pure rational admits something better. If \(p/q\) is reported as \(v\) at \(N\) places then

\[
  v\,q \;\le\; p \cdot 10^{N} \;<\; (v+1)\,q
\]

is an exact integer statement with no arithmetic in it that could go wrong. Either it holds or the
digits are wrong. That is \textbf{verification}, and not agreement, and it is available for the
thirty eight constants that do not involve \(\pi\).

The seven defining constants are integers times powers of ten:

\begin{center}
\begin{tabular}{lr}
\toprule
constant & exact value \\
\midrule
caesium hyperfine frequency & \(9\,192\,631\,770\) Hz \\
speed of light & \(299\,792\,458\) m/s \\
Planck constant & \(662607015 \times 10^{-42}\) J s \\
elementary charge & \(1602176634 \times 10^{-28}\) C \\
Boltzmann constant & \(1380649 \times 10^{-29}\) J/K \\
Avogadro constant & \(602214076 \times 10^{15}\) mol\(^{-1}\) \\
luminous efficacy & \(683\) lm/W \\
\bottomrule
\end{tabular}
\end{center}

Everything built from those by field operations is a rational. This can be
exact at all for that reason.

\section{The error the exact check could not catch}

\texttt{elementary charge over h-bar} came out wrong on the first run, and how it was caught is the
argument for keeping an external check even when an internal one is exact.

\[
  e/\hbar \;=\; e \big/ (h/2\pi) \;=\; 2\pi e / h
\]

Every other \(\pi\) constant in the set has \(\hbar\) \emph{as} its value and therefore divides by
\(2\pi\). This one has \(\hbar\) in its \textbf{denominator}, and multiplies. A division here,
with a spurious factor of two besides, would come out low by
\((2\pi)^2 = 19.74\).

\textbf{The exact check would pass it.} A multiplication back verifies a correct division of a rational,
and says nothing about whether it is the right rational. An exact check cannot see a correctly
computed wrong expression. The NIST comparison catches it.

That has the same shape as the standing lesson on instruments in this tree: a blindness is a property
of the instrument, and it must be known, and not hoped about.
